% a01_H0rel part 3: face-localised differences, the constant datum, the derivative clause for small-field data
\subsection{The face-localised difference and the background of small-field data}\label{app:a01:fld}

\begin{lemma}[reproduction of the constant]\label{lem:a01:RC}
On the doubled torus $2X$ with the periodic extension of the weight (the killed sites of the cube being ordinary zero-weight sites there), and on the cube, $a_jG_jQ^*\mathbf 1=\mathbf 1-V_jG_j\mathbf 1$ at $A=0$, where on the cube $G_j\defeq\Mop(\Xl,0)^{-1}$ and the identity holds up to the face term of \cref{lem:a01:POI}; and $\norm{G_j\mathbf 1}_\infty\le C$.
\end{lemma}
\begin{proof}
$Q\mathbf 1=\mathbf 1$ (Q1), so $P\mathbf 1=Q^*\mathbf 1$. With $v\defeq\mathbf 1-V_jG_j\mathbf 1$: $(-\Delta+V_j+a_jP)v=V_j+a_jP\mathbf 1-V_j=a_jQ^*\mathbf 1$ on the torus, where $-\Delta\mathbf 1=0$. On the cube $-\Delta^D\mathbf 1=\mathsf B$, the face term of \cref{lem:a01:POI}. The sup bound is \cref{lem:a01:A4}(c).
\end{proof}

\begin{lemma}[the Dirichlet Poisson term decays]\label{lem:a01:POI}
Let $\mathsf B(x)\defeq\vth^{-2}\#\{\text{killed neighbours of }x\}$ (supported on $F_1$) and $v_0\defeq\Mop(\Xl,0)^{-1}\mathsf B$. Then $|v_0(x)|\le Ce^{-(\delta_2/2)(t(x)-4)}$ for every $x$, and $|D_0v_0(b')|\le Ce^{-(\delta_2/2)(t(b')-5)}$ at every bond with $t\ge6$.
\end{lemma}
\begin{proof}
(1) $\Mop(\Xl,0)\mathbf 1=\mathsf B+V_j+a_jQ^*\mathbf 1$ on $\Xl$, so $h^0\defeq a_j\Mop^{-1}Q^*\mathbf 1=\mathbf 1-\Mop^{-1}V_j-v_0$ and $v_0$ is bounded by $1+Cc_V+C$ (\cref{lem:a01:A4}). (2) Let $\chi(x)\defeq\prod_\mu\psi(x_\mu)\psi(N-x_\mu)$ with $\psi$ smooth, $0$ on $(-\infty,3.25]$, $1$ on $[3.75,\infty)$, constant on a $\vth$-neighbourhood of every block plane. Then $\Mop(\chi v_0)=[\Mop,\chi]v_0$ (since $\mathsf B=0$ where $\chi\ne0$), supported on the blocks at face distance in $[3,4]$; the Caccioppoli inequality on a two-block neighbourhood (where $\Mop v_0=0$; the bond sum over $\Xl$ has no boundary term, the face bonds entering with $v_0=0$ at $\Kill$) bounds $\norm{\mathbf 1_B[\Mop,\chi]v_0}_2\le C\sup|v_0|\le C$, and \cref{lem:a01:A4}(c) with block-$L^2$ data gives $|\chi v_0(x)|\le C\sum_{B\in\text{shell}}e^{-\delta_2\dist(x,B)}\le Ce^{-(\delta_2/2)(t(x)-4)}$. (3) The interior estimate on $B_{1/4}(x)$ at $t\ge6$ with $g=-V_jv_0-a_jPv_0$.
\end{proof}

\begin{lemma}[face-localised difference, whole cube, regular $A$]\label{lem:a01:FLD}
Let $T\defeq2X$ carry the periodic extension of the weight and the even extension $A_T\defeq\Ebar A$ of a regular field $A$ (regular with the same constant by \cref{lem:a01:REG}(b)), and $\Mop_T\defeq-\Delta_{A_T}+V_j+a_jP^T$. For bounded $F$ supported in $\Xl$ put $v\defeq\Mop_T^{-1}F$, $u\defeq\Mop(\Xl,A)^{-1}F$ and $w\defeq v|_{\Xl}-u$. Then for every site $x$ and fine bond $b$ of the cube,
\[
|w(x)|+|D_Aw(b)|\ \le\ C\exp\big[-\delta_F\big(\dist(\cdot,\DD)+\dist(\supp F,\DD)+\dist(\cdot,\supp F)\big)\big]\,\norm F_\infty,
\]
with $C,\delta_F$ depending on $(d,L,a,\Wt,c_V)$ and King's constants only.
\end{lemma}
\begin{proof}
(1) For $x\in\Xl$, $(\Mop_Tv)(x)$ and $(\Mop(\Xl,A)v|_X)(x)$ differ only by the hopping to killed neighbours, since $P^T(x,\cdot)$ is supported in the block of $x$, whose weighted sites are live (Q8). So $(\Mop(\Xl,A)v|_X)(x)=F(x)+h(x)$ with $h(x)\defeq\vth^{-2}\sum_{x'\in\Kill,x'\sim x}U_{xx'}v(x')$ supported on $F_1$, and $w=\Mop^{-1}h$. (2) Near the face, $\dist(x,\DD)\le D_0\defeq4+2\sqrt3$: $|w|\le|v|+|u|$ and $|D_Aw|\le|D_Av|+|D_Au|$; on the torus $|v(x)|\le Ce^{-\delta\dist(x,\supp F)}\norm F$ (\cref{lem:a01:A4}(c) on $T$) and $|D_Av(b)|\le Ce^{-\delta\dist(b,\supp F)/2}\norm F$ (interior estimate on unit balls); on the cube the same for $u$ up to the face (\cref{lem:a01:A4}(c), \cref{lem:a01:A5} in the bounded-source form). As $\dist(\supp F,\DD)\le\dist(x,\supp F)+D_0$ near the face, the claim holds there with $\delta_F\le\delta/4$. (3) Away from the face: let $\chi$ be smooth, $\chi=0$ on $\{\dist(\cdot,\DD)\le1+\sqrt3\}$, $\chi=1$ on $\{\dist\ge D_0-1\}$, $|\nabla\chi|\le1$, $|\nabla^2\chi|\le C$. Then $\chi h=0$ and $\Mop(\chi w)=[\Mop,\chi]w=:g$ with $g=[-\Delta_A,\chi]w+a_j[P,\chi]w$; the lattice Leibniz rule gives $|[-\Delta_A,\chi]w|\le C(|w|+|D_Aw|)$ on the shell $S\defeq\{\sqrt3\le\dist(\cdot,\DD)\le D_0\}$, and $|[P,\chi]w(x)|\le\sum_{x'}\vth^3|P(x,x')||\chi(x)-\chi(x')||w(x')|\le2\sqrt3\,\Wt\sup_{B(x)}|w|$ by (Q6), (Q8), supported on the blocks meeting the shell. So $g$ is bounded, supported in $S'$ (the union of $S$ with the blocks meeting it, at distance $\le D_0+\sqrt3$ from $\DD$ and $\ge1$ from $\Kill$), and by (2) $|g(z)|\le Ce^{-\delta\dist(z,\supp F)/2}\norm F$. (4) By \cref{lem:a01:A4}(c) in the decaying form, $|\chi w(x)|\le C\sup_{z\in S'}e^{-\delta_2|x-z|/2}|g(z)|$, hence for $\dist(x,\DD)\ge D_0$: $|w(x)|\le C\sup_{z\in S'}e^{-\delta_2|x-z|/2-\delta\dist(z,\supp F)/2}\norm F$; with $|x-z|\ge\dist(x,\DD)-D_0-\sqrt3$, $\dist(z,\supp F)\ge\dist(\supp F,\DD)-D_0-\sqrt3$ and $|x-z|+\dist(z,\supp F)\ge\dist(x,\supp F)$, averaging the exponents gives the claim for $|w|$ with $\delta_F\defeq\min(\delta,\delta_2)/6$. (5) For $\dist(x,\DD)\ge D_0$, $w$ solves $-\Delta_Aw=-V_jw-a_jPw$ on $B_{1/4}(x)$ ($h=0$ there), so the interior estimate gives $|D_Aw|\le C\Wt\sup_{B_{1/2}(x)+\text{one block}}|w|$, which obeys the bound of (4) with $D_0$ replaced by $D_0+1$.
\end{proof}

The comparison operator is the torus operator with an arbitrary regular extension of $A$; the scalar is never reflected, so no regularity of a reflected field is needed. No flat-face condition is used: the cut-off sees only $\dist(\cdot,\DD)$, so edges and corners of the cube are included.

\begin{lemma}[the Dirichlet image effect on a sub-domain, sup form]\label{lem:a01:LOC}
Let $\Omega$ be a sub-domain and $\Omega^\sharp$ a Neumann domain without Dirichlet face that contains $\Omega\cap\Xl$ and the killed sites adjacent to it, adds no bond from $\Omega\cap\Xl$ except the bonds to those sites, and carries a regular extension $A^\sharp$ of $A$ (for instance $\Omega$ together with its mirror blocks across the Dirichlet planes it meets, with the even extension of every component of $A$). For data $b$ on the blocks of $\Omega$ (zero on the added blocks) let $u\defeq a_j\Mop(\Omega,A)^{-1}Q^*b$ and $u^\sharp\defeq a_j\Mop(\Omega^\sharp,A^\sharp)^{-1}Q^*b$. Then for every $x\in\Omega\cap\Xl$, $|u(x)-u^\sharp(x)|\le C_U|b|_\infty e^{-(\delta_2/2)(t(x)-4)}$ with $C_U=C_U(d,a,\Wt,\gamma)$, free of Ba{\l}aban's large-block size $M$, of $\vth$, $j$, $\Omega$ and $A$; and $|D_A(u-u^\sharp)(b')|\le C|b|_\infty e^{-(\delta_2/2)(t(b')-5)}$ at bonds with $t\ge6$ and depth $\ge1$.
\end{lemma}
\begin{proof}
$w\defeq u-u^\sharp|_\Omega$ solves $\Mop(\Omega)w=0$ at every site of $\Omega\cap\Xl$ off $F_1$ (same bonds, same blocks; the killed sites have zero weight). The product cut-off of \cref{lem:a01:POI}, the Caccioppoli inequality in the region $t\ge2.25$, $\sup|w|\le|u|+|u^\sharp|\le C|b|$ (\cref{lem:a01:A4} on both domains) and \cref{lem:a01:A4}(c) with block-$L^2$ data give the sup bound exactly as in \cref{lem:a01:POI}; the derivative bound is the interior estimate.
\end{proof}

\begin{lemma}[transported constant up to the face]\label{lem:a01:H}
Let $A$ be regular on $\Xl$, $b$ a fine bond with base $x_0$, $\lambda_0(z)\defeq\sum_\mu A_\mu(x_0)z_\mu$, $A'\defeq A-d\lambda_0$ (gauge equivalent to $A$ on the simply connected cube, with $A'=0$ on the bonds at $x_0$ and $|A'_{b''}|\le K_j|b''-x_0|_1$), and $h\defeq a_j\Mop(\Xl,A')^{-1}Q(A')^*\mathbf 1$. Then $|D_{A'}h(b)|\le C_H(V_j+\theta_j+e^{-\delta_ht(x_0)})$ with $\delta_h\defeq\delta_2/2$ and $C_H=C_H(d,L,a,\Wt,\gamma,c_V)$.
\end{lemma}
\begin{proof}
If $t(x_0)<6$ the claim is \cref{lem:a01:A5} ($h$ is the leg of the datum $\mathbf 1$). Let $t(x_0)\ge6$. (1) At $A=0$: by \cref{lem:a01:POI}(1), $D_0h^0=-D_0\Mop_0^{-1}V_j-D_0v_0$, so $|D_0h^0(b)|\le CV_j+Ce^{-\delta_h(t-5)}$ (\cref{lem:a01:A5} for the bounded source $V_j$, \cref{lem:a01:POI}). (2) $w\defeq h-h^0=\Mop(A')^{-1}S$ with $S\defeq a_j(Q(A')^*-Q(0)^*)\mathbf 1-(\Mop(A')-\Mop(0))h^0$. Pointwise $|S(z)|\le C\theta_j(1+|z-x_0|)^2$: the block terms by $|U(A'(\Gamma))-1|\le e_j|A'||\Gamma|\le C\theta_j(1+|z-x_0|)$ and $|h^0|\le C$; the Laplacian terms by pairing the two bonds of each direction, $\vth^{-2}[(1-U_+)+(1-U_-)]h^0(z)+\vth^{-2}\sum_\pm(1-U_\pm)(h^0(z_\pm)-h^0(z))$, where the first bracket is $-\vth e_jq(a_+-a_-)+O(\vth^2e_j^2|A'|^2)$, so its $\vth^{-2}$-multiple is $\le\theta_j+C\theta_j^2|z-x_0|^2$, and the second is $\le\sum e_j|A'||D_0h^0|\le C\theta_j|z-x_0|$ (\cref{lem:a01:A5} at $A=0$); at a first-layer site the normal pair is incomplete and the unpaired term is $\le\vth^{-1}e_j|A'|\cdot2C_{\mathrm{DV}}\vth\le C\theta_j|z-x_0|$ by \cref{lem:a01:DV}. (3) $|w(x)|\le C\theta_j(1+|x-x_0|)^2$ by \cref{lem:a01:A4}(c), shell by shell around $x_0$. (4) The interior estimate for $w$ on $B_{1/4}(x_0)$ (no killed site, $t\ge6$) with $-\Delta_{A'}w=S-V_jw-a_jP(A')w$, bounded by $C\theta_j$ on $B_{1/2}(x_0)$, gives $|D_{A'}w(b)|\le C\theta_j$. (5) $D_{A'}h(b)=D_0h^0(b)+(D_{A'}-D_0)h^0(b)+D_{A'}w(b)$, and the middle term vanishes since $A'=0$ on the bonds at $x_0$.
\end{proof}

\begin{proposition}[derivative clause for backgrounds of small-field data]\label{prop:a01:legsd}
Let $A$ be regular on $\Xl$, $\phi_j$ block data, $b$ a fine bond with base $x_0$ in the block $y_0$, $r_j\defeq R(1+\ln\ell_j^{-1})^{\varrho}$ Ba{\l}aban's region radius, and assume: (G) $|D_{A_j}\phi_j(\beta)|\le p_j$ on every block bond within block distance $2r_j$ of $y_0$; (S) $|\phi_j|\le M_j$; (F) King's face test $|\phi_j(y)|\le\sigma_j\le C_8p_j$ on the face blocks within $2r_j$ of $y_0$; and $M_j(V_j+\theta_j)\le p_j$, $M_je^{-\delta_*r_j}\le p_j$ with $\delta_*\defeq\min(\delta_3,\delta_h)/2$ (both hold for $\eta\le\eta_1$ and $R\ge R_*$: $M_j\theta_j\le Ce\lambda^{-1/4}L^{1/4}(1+\ln L)^pp_j^2\ell_j^{1/4}$). Then the background $\phi^{(j)}=a_j\Mop(\Xl,A)^{-1}Q(A)^*\phi_j$ satisfies $|D_A\phi^{(j)}(b)|\le C_{\mathrm{legs}}p_j$ at every fine bond of the cube, face bonds included, $C_{\mathrm{legs}}=C(1+C_8)$. On a sub-domain $\Omega$, under the same hypotheses for the data on $\Xl$, the domain background $\phi^{(j)}_\Omega\defeq a_j\Mop(\Omega,A)^{-1}Q^*(\phi_j|_\Omega)$ satisfies, at every bond of depth $\ge\frac12$,
\[
|D_A\phi^{(j)}_\Omega(b)|\ \le\ C_{\mathrm{legs}}p_j+CM_je^{-\delta'(\dep(b)-4)_+},\qquad\delta'\defeq\delta_3/2,
\]
hence $\le(C_{\mathrm{legs}}+1)p_j$ at depth $\ge\rho_j\defeq4+\delta'^{-1}\ln(CM_j/p_j)$, and $\rho_j\le r_j$ once $R\ge R'_*$.
\end{proposition}
\begin{proof}
Gauge to $A'$ as in \cref{lem:a01:H} (every quantity is gauge covariant). (1) For blocks $y$ with $|y-y_0|\le r_j$ let $U_y$ be the coarse transporter along the coordinate-ordered block path from $y_0$ to $y$ (inside the convex cube); telescoping (G) gives $\phi_j(y)=U_yc+\psi(y)$, $c\defeq\phi_j(y_0)$, $|\psi(y)|\le p_j|y-y_0|_1$; for $|y-y_0|>r_j$ put $\psi(y)\defeq\phi_j(y)$, $|\psi|\le M_j$. In the gauge $A'$, $|U_y-1|\le C\theta_j(1+|y-y_0|)^2$. (2) The $\psi$-part: $|D_{A'}[a_j\Mop^{-1}Q^*\psi](b)|\le\sum_y|D_{A'}l(b,y)||\psi(y)|\le C(p_j+M_je^{-\delta_3r_j})\le Cp_j$ by the decaying clause of \cref{lem:a01:A5}; the transport correction $|c||D_{A'}[a_j\Mop^{-1}Q^*((U-1)\mathbf 1_{|y-y_0|\le r_j})](b)|\le CM_j\theta_j\le Cp_j$; the truncation $|c||D_{A'}[a_j\Mop^{-1}Q^*\mathbf 1_{|y-y_0|>r_j}](b)|\le CM_je^{-\delta_3r_j}\le Cp_j$. (3) The constant part: $|c||D_{A'}h(b)|\le|c|C_H(V_j+\theta_j+e^{-\delta_ht(x_0)})$ by \cref{lem:a01:H}. (4) The cap: if $t(x_0)\le r_j$, transport along the normal block path from $y_0$ to the face layer (length $\le t(x_0)+1$) and (F) give $|c|\le\sigma_j+p_j(t(x_0)+1)$, so $|c|e^{-\delta_ht}\le C(1+C_8)p_j$; if $t>r_j$, $|c|e^{-\delta_ht}\le M_je^{-\delta_hr_j}\le p_j$; and $|c|C_H(V_j+\theta_j)\le Cp_j$. On $\Omega$: $w\defeq\phi^{(j)}_\Omega-\phi^{(j)}|_\Omega$ solves $\Mop(\Omega)w=0$ at depth $\ge1$ (same bonds and blocks there); with $\chi\defeq\psi(\tilde d_N)$, $\tilde d_N$ a mollified depth constant near block planes, $\Mop(\Omega)(\chi w)=[\Mop,\chi]w=:F$, supported at depth in $[3.25,3.75]$ plus one block, with $|F|\le C(|w|+|D_Aw|+\sup_{\text{block}}|w|)\le CM_j$ (\cref{lem:a01:A4} for both backgrounds, \cref{lem:a01:A5} at depth $\ge3$); \cref{lem:a01:A5} for bounded sources gives $|D_A(\chi w)(b)|\le CM_je^{-\delta_3(\dep(b)-4)}$, and $\chi=1$ near $b$ when $\dep(b)\ge4$; for $\frac12\le\dep(b)<4$ use \cref{lem:a01:A5} on both backgrounds.
\end{proof}

The face test is necessary: for the constant data $\phi_j\equiv c$ the background vanishes at the Dirichlet face and $|D_A\phi^{(j)}|$ is of order $|c|$ on the first block; this is why King imposes \cite[(5.41)]{K2} at each step.

\begin{lemma}[face-localised background of a covariantly constant datum, regular $A$]\label{lem:a01:FL}
Let $\Omega$ be a box of large blocks with Dirichlet conditions on its faces (the cube, or a sub-box), $y_0$ a block, $c\in\C$, $\mathsf T$ a tree of coordinate-ordered block paths from $y_0$, $\Theta_A(y)\defeq\overline{U(A(\Gamma_{\mathsf T}(y_0\to y)))}\,c$ the datum covariantly constant along $\mathsf T$, and $u_A\defeq a_j\Mop(\Omega,A)^{-1}Q(A)^*\Theta_A$. There are $\delta_1>0$, $C$ and $\theta_*>0$ such that for $A$ regular on $\Omega$ with $\theta_j\le\theta_*$, every bond $b$ of $\bar\Omega$ (face bonds included) and every block $y$: (a) $|D_{A,b}u_A|\le|c|[CV_j+Ce^{-\delta_1\dist(b,\DD)}+C\theta_j(1+|b-y_0|)^3]$; (b) $|\Theta_A(y)-(Q(A)u_A)(y)|\le|c|[CV_j+Ce^{-\delta_1\dist(B(y),\DD)}+C\theta_j(1+|y-y_0|)^3]$; (c) $\sup|u_A|\le C|c|$.
\end{lemma}
\begin{proof}
Step 1 ($A=0$, $\Theta=c$). The Dirichlet operator of a box of large blocks with King's weight is the odd restriction of the free operator on $\vth\Z^3$ with the periodic weight: the odd extension of a function vanishes on the Dirichlet planes, the weight couples no live site across a plane, and the free operator preserves the odd subspace, on which it is unitarily $\Mop(\Omega,0)$ (coercivity gives uniqueness). So $\Mop(\Omega,0)^{-1}F=(G^{\mathrm{free}}\Ebar_{\mathrm{odd}}F)|_{\bar\Omega}$ and $\Ebar_{\mathrm{odd}}Q^*c=Q^*(cs)$ with $s(y)=\pm1$ the coarse odd extension. Write $s=1-(1-s)$: by \cref{lem:a01:RC} on the free lattice, $a_jG^{\mathrm{free}}Q^*\mathbf 1=\mathbf 1-V_jG^{\mathrm{free}}\mathbf 1$, so $u_0=c[\mathbf 1-V_jG^{\mathrm{free}}\mathbf 1-a_jG^{\mathrm{free}}Q^*(1-s)]$ on $\bar\Omega$; $(1-s)$ is supported on the image boxes with $s=-1$, at distance $\ge\dist(b,\DD)$ from $b$, and bounded by $2$. With the weighted-$\ell^1$ bounds $\norm{DG^{\mathrm{free}}}_{\delta_1},\norm{G^{\mathrm{free}}}_{\delta_1}\le C$ of the free weighted operator (\cref{lem:a01:A23,lem:a01:A4} on $\vth\Z^3$), $|Du_0(b)|\le|c|[V_j\norm{DG^{\mathrm{free}}}_0+2a_j\Wt\norm{DG^{\mathrm{free}}}_{\delta_1}e^{-\delta_1\dist(b,\DD)}]$ at every bond, face bonds included (the intertwining holds there); the same computation with $Q$ in place of $D$ gives (b), and (c) is $\norm{G^{\mathrm{free}}}_0$. Step 2 (constant $B$): gauge covariance, $u_B=U_\lambda u_0(c')$ with $|c'|=|c|$, $\lambda=B\cdot x$; (a)--(c) hold with the $\theta$ term absent. Step 3 (regular $A$). Fix the bond $b$ and let $B$ be the constant field with the components of $A$ at reference bonds through $b$; $A'\defeq A-B$ obeys $e_j|A'(x)|\le d\theta_j(1+|x-b|_1)$ on the box (sum the first differences along a lattice path). With $w\defeq u_A-u_B$ (both vanish on $\DD$): $\Mop_Aw=F\defeq a_j[Q(A)^*\Theta_A-Q(B)^*\Theta_B]+(\Mop_B-\Mop_A)u_B$; the hopping difference is a first-order operator with coefficients $\le C\theta_j(1+|x-b|)$ and the block part has row sums $\le C\theta_j(1+|x-b|)$, the contour transports differ by $\le e_j\sum_\Gamma\vth|A'|$, and $|u_B|,|D_Bu_B|\le C|c|$ by Steps 1--2; so $|F(x)|\le C\theta_j|c|(1+|x-b|+|y_0-b|)^3$. Split $F=\sum_nF_n$ by shells $n\le|x-b|<n+1$; \cref{lem:a01:A5} in the bounded-source form (face bonds included) gives $|D_{A,b}\Mop_A^{-1}F_n|\le Ce^{-\delta'(n-1)}\norm{F_n}_\infty$ and the same for $|\Mop_A^{-1}F_n|$ within one block of $b$; summing, $|D_{A,b}w|+\sup_{|x-b|\le2}|w|\le C\theta_j|c|(1+|y_0-b|)^3$. Finally $D_{A,b}u_B=D_{B,b}u_B$ since $A$ and $B$ agree on $b$ (on a face bond both are the gauge-free $\vth^{-1}|u_B|$), which gives (a); (b) from the four-term split $\Theta_A-Q(A)u_A=[\Theta_B-Q(B)u_B]+[\Theta_A-\Theta_B]-[(Q(A)-Q(B))u_B]-Q(A)w$, and (c) from the sup clause with $\norm F_\infty=a_j\Wt|c|$.
\end{proof}

The face rate $\delta_1$ is the free rate of the weighted operator at $A=0$ and does not depend on Ba{\l}aban's large-block size $M$; this is what allows the near-layer constant $R_0$ of \cref{app:a01:upper} to be fixed before $M$.

\begin{lemma}[the block-mismatch identities]\label{lem:a01:EF}
\label{lem:a01:EB}
For every weight, domain, boundary condition and gauge field, with $\phi^{(j)}=a_j\Mop^{-1}Q(A)^*\phi_j$: (a) $\phi_j-Q(A)\phi^{(j)}=a_j^{-1}\Delta^{(j)}(A)\phi_j$; (b) with the \emph{uniform} transported average $Q^u(A)\psi(y)\defeq\sum_{x\in B(y)}\vth^3U(A(\Gamma_{y,x}))\psi(x)$, exactly $a_j(\phi_j-Q(A)\phi^{(j)})(y)=Q^u(A)[(-\Delta_A+V_j)\phi^{(j)}](y)$, whereas with $Q(A)$ in place of $Q^u(A)$ the left side acquires the factor $W_2$.
\end{lemma}
\begin{proof}
(a) $Q\phi^{(j)}=a_jQ\Mop^{-1}Q^*\phi_j$ and $\Delta^{(j)}=a_j-a_j^2Q\Mop^{-1}Q^*$. (b) The minimiser equation is $(-\Delta_A+V_j)\phi^{(j)}=a_jQ^*(\phi_j-Q\phi^{(j)})=:a_jQ^*m$, and pointwise $(Q^*m)(x)=w(x)U(\Gamma_{y,x})^*m(y)$ for $x\in B(y)$; applying $Q^u(A)$ gives $a_jm(y)\sum_x\vth^3w(x)=a_jm(y)$ by (W1) and $UU^*=1$; applying $Q(A)$ gives $a_jW_2m(y)$.
\end{proof}

Identity (b) is the form in which the block-mismatch clause of the lower bound is read on the cube; the naive summation by parts of the \emph{weighted} average $Q[-\Delta\phi](y)$ produces, besides the boundary flux of the block, interior terms $\sum_{\text{bonds}}\vth^{d-2}(w(x)-w(x'))(\phi(x)-\phi(x'))$ across the zero-weight layers, whose absolute bound $2d\Wt\sup|D\phi|/(L\vth)$ diverges as $\vth\to0$; the uniform transported average has no interior terms, and (a) needs no fine-level summation by parts at all.
